
Scene-graph datasets for robotics are built by hand: every spatial relationship between every pair of objects is decided by a human annotator. This bottleneck keeps such datasets small, and the manual labels it produces are inconsistent: problems the authors of a recent robot-acquired spatial-relationship dataset report themselves. This dissertation asks whether the human can be removed from the labelling loop entirely, for the seven spatial predicates the dataset defines (on, under, left/right of, in front of/behind, near).

The proposed pipeline computes, rather than predicts, each relationship: off-the-shelf perception models measure the scene (segmentation masks and monocular depth), each object is lifted to a 3D position, and explicit geometric rules decide every label. Rule thresholds are fitted only on six of the dataset's nine annotator groups and validated on the held-out three, and outputs are written in the dataset's native formats. The full 836-image dataset is annotated in about five minutes on a consumer GPU, with 20 times the label density of the manual pass.

Validated against 8,926 human-annotated relationships, the automatic labels match or exceed the human process on five of seven predicates (0.85 mean recall, 0.76 on held-out annotators; audited precision $\approx$ 1.0 for the lateral and proximity predicates and $\approx$ 0.9 for support). The hardest pair, in front of/behind, is decided by a two-stage cascade of relative depth and a ground-plane projection cue, and diagnosing every remaining disagreement attributes the residual gap to measured properties of the human annotation itself, including two annotator groups that labelled the pair with opposite conventions, rather than to tool error (\textasciitilde{}7\% of misses). In a controlled downstream experiment, a classifier trained on the automatic labels reaches 0.76 mean recall against held-out human annotations, versus 0.30 when trained on the human labels and 0.36 when those labels are stretched by self-training, the standard semi-supervised remedy: at this dataset's annotation scale, dense and consistent computed labels are better training material than the sparse human labels they replace, and better than any attempt to extrapolate from them.

Repeating the comparison in the source paper's own benchmark framework, with a shared frozen detector and three seeds per arm, splits the verdict. Trained on human labels the model ranks better against the human test annotation (mR@100 0.326 against 0.278), but on relation compositions never seen during training it recalls almost nothing, while the automatically trained model recalls sixty times more (zR@100 0.172 against 0.003). Scored per annotator, the human-trained advantage is large against both annotators carrying a measured labelling defect and disappears against the only one without. The ranking metric therefore rewards agreement with annotation habits as well as spatial correctness, which is a property of the benchmark rather than of either label source, and the dissertation's answer is correspondingly conditional: automatic labels are the better supervision wherever ground truth means geometry, human labels wherever it means annotation practice. Robot planning needs the former. The annotation bottleneck this pipeline removes was not only limiting dataset size; it was limiting what the dataset could teach.